%% SORI ↔ SoundPLAN 계산 조건 대조표 (학술발표용). 단독 컴파일 문서: xelatex 2회
%% 근거: 소음진동설계/SoundPLAN_계산조건_정리.pdf(2026-11-17) 4절 표 「앱별로 SoundPLAN과 다른 설정」 + 1절 설정표
%%       SORI 설정 범위: sori/config.py(DEFAULTS·_validate), sori/geometry.py(MAXO), sori/rls19.py(표 4a·4b·5)
%%       SORI 이번 계산 설정: sp_refl/config.yaml, sp_refl_1119/config.yaml(교통량만 SoundPLAN 입력값으로 바꿈)
%% 본 문서는 Claude(Anthropic의 AI)를 활용하여 작성하였음.
\documentclass[10pt,a4paper]{article}
\usepackage[a4paper,top=15mm,bottom=15mm,left=16mm,right=16mm]{geometry}
\usepackage{fontspec}
\usepackage{xetexko}
\IfFontExistsTF{Noto Serif CJK KR}{%
  \setmainfont{Noto Serif CJK KR}\setsansfont{Noto Sans CJK KR}%
  \setmainhangulfont{Noto Serif CJK KR}\setsanshangulfont{Noto Sans CJK KR}%
  \setmonohangulfont{Noto Sans CJK KR}[Scale=0.85]}{%
  \setmainfont{Times New Roman}\setsansfont{Apple SD Gothic Neo}%
  \setmainhangulfont{AppleMyungjo}[AutoFakeBold=2.5]\setsanshangulfont{Apple SD Gothic Neo}%
  \setmonohangulfont{Apple SD Gothic Neo}[Scale=0.85]}
\usepackage{booktabs,tabularx,array,longtable,xcolor,colortbl}
\usepackage{enumitem,needspace}
\usepackage[hidelinks]{hyperref}
\linespread{1.12}
\setlength{\parindent}{0pt}
\setlength{\parskip}{0.35em}
\setlist{itemsep=0.1em,topsep=0.25em}
\renewcommand{\tablename}{표}
\newcolumntype{L}{>{\raggedright\arraybackslash}X}
\newcolumntype{P}[1]{>{\raggedright\arraybackslash}p{#1}}
\newcolumntype{C}{>{\centering\arraybackslash}X}
\newcommand{\f}[1]{\texttt{\small #1}}
\newcommand{\hl}[1]{\textbf{#1}}
\definecolor{okbg}{HTML}{E3F2E6}\definecolor{okfg}{HTML}{1E6B33}
\definecolor{nbbg}{HTML}{FDEBD8}\definecolor{nbfg}{HTML}{A4500F}
\definecolor{hdr}{HTML}{DCE4F2}\definecolor{navy}{HTML}{14285A}
\newcommand{\ok}{\colorbox{okbg}{\textcolor{okfg}{\sffamily\bfseries\footnotesize 같음}}}
\newcommand{\eq}[1]{\colorbox{okbg}{\textcolor{okfg}{\sffamily\bfseries\footnotesize #1}}}
\newcommand{\nb}[1]{\colorbox{nbbg}{\textcolor{nbfg}{\sffamily\bfseries\footnotesize #1}}}
\newcommand{\sect}[1]{\Needspace{6\baselineskip}\section*{\textcolor{navy}{#1}}}

\begin{document}

{\sffamily\Large\bfseries\textcolor{navy}{SORI 계산 조건 — SoundPLAN 계산 조건과의 대조}}\\[1mm]
{\sffamily 전농사거리 사가정로 모델 · SoundPLAN 9.1 / SORI v0.9 · RLS-19 · 2026년 11월 19일 · 전현수}

\vspace{2mm}
\fcolorbox{blue!50!black}{blue!5}{\parbox{\dimexpr\linewidth-2\fboxsep-2\fboxrule}{\small
\textbf{요약.} 「SoundPLAN 계산 조건 정리」(2026-11-17) 4절 표에서 SORI 열은 「설정 가능 · 기본값 가정 · 확인」으로 비어 있었다.
이 문서는 그 칸을 채우고, 항목마다 SORI의 계산 조건(이번 비교에 넣은 값)·기본값·변경 가능 범위를 정리한다.
\begin{itemize}[leftmargin=1.2em]
\item SORI는 배출·전파를 \hl{RLS-19}로 계산하도록 설정했고, 4절 표의 7개 항목과 표 아래 4개 확인 항목(반사 손실, 반사 거리 제한, 차폐량 상한, LrD)을
  \hl{모두 SoundPLAN과 같게} 맞췄다.
\item 1절 설정표의 나머지 항목(도로 분할, 탐색 반경, 교통량·속도, 노면, 도로 높이, 건물·지형, 수음점)도 같게 맞췄다.
\item 같게 맞출 수 없는 항목은 \hl{격자 보간}(SoundPLAN 9×9 보간, SORI는 모든 점을 직접 계산) 하나다.
  해석이 들어간 항목은 \hl{반사면–음원 50 m 판정 방식} 하나다(4절).
\end{itemize}}}

%% ------------------------------------------------------------------
\sect{1. 「앱별로 SoundPLAN과 다른 설정」 표의 SORI 열}
「SoundPLAN 계산 조건 정리」 4절 표의 항목 순서를 그대로 따랐다. SORI 칸은 세 가지로 나눴다.
\hl{계산 조건}은 SoundPLAN과 비교하려고 이번에 넣은 값, \hl{기본값}은 따로 지정하지 않을 때 쓰는 값, \hl{변경 가능 범위}는 설정으로 바꿀 수 있는 값이다.

{\footnotesize
\renewcommand{\arraystretch}{1.25}
\begin{tabularx}{\linewidth}{@{}>{\bfseries}P{19mm}P{22mm}P{28mm}P{30mm}Lc@{}}
\toprule
\rowcolor{hdr} 항목 & SoundPLAN & SORI 계산 조건 & SORI 기본값 & SORI 변경 가능 범위 & 판정 \\
\midrule
반사 차수 & 2차 & \hl{2차} & 3차(「소음지도의 작성방법」 별표 2 「3차 이상」) &
  \hl{0 \textasciitilde{} 6차}. 격자·외벽·점 수음점마다 따로 지정. 4차 이상은 계산 시간이 차수마다 크게 늘어남 & \ok \\
노면 보정 & 0 dB (기준 노면) & \hl{기준 노면}(D\textsubscript{SD} = 0 dB) & 아스팔트·콘크리트 → 기준 노면(0 dB), 블록·비포장 → 블록 포장 &
  RLS-19 표 4a 노면 11종(기준, SMA, AC 11, 배수성 PA 8·11 등) + 표 4b 블록 포장 2종. 수치지도 「포장재질」 값마다 지정 & \ok \\
신호교차로 보정 & 없음 & \hl{끔}(D\textsubscript{K} = 0) & 켬: 4차로 이상 도로 2개가 만나는 교차점을 신호교차로로 자동 판정 &
  켬/끔. 자동 판정 기준 차로 수 또는 교차점 CSV 직접 지정. RLS-19 표 5: 신호 3 dB, 회전 2 dB & \ok \\
다중반사 D\textsubscript{refl} & 없음 & \hl{끔}(D\textsubscript{refl} = 0) & 켬(도로 양쪽 건물 자동 탐색) &
  켬/끔. RLS-19 식 (8) & \ok \\
건물 높이 & 층수 × 3 m & \hl{층수 × 3 m} & 층수 × 3 m &
  층고 임의 값, 또는 실제 높이 필드 사용. 무벽건물·4 m\textsuperscript{2} 미만 건물 제외 & \ok \\
측면 회절 & 없음 & \hl{없음}(RLS-19 전파) & 켬(ISO 9613-2 전파일 때만 적용) &
  ISO 9613-2 전파에서 켬/끔(식 13, 건물 감싸 돌기 0 \textasciitilde{} 20회). RLS-19·RLS-90 전파에서는 원문 규정대로 항상 없음 & \ok \\
기준 표기 & RLS-19 & \hl{RLS-19}(배출·전파) & RLS-19 배출 + ISO 9613-2 전파 &
  RLS-19(배출·전파) / RLS-19 배출 + ISO 9613-2 전파 / RLS-90(배출·전파) & \ok \\
\midrule
\multicolumn{6}{@{}l}{\sffamily 4절 표 아래 「모든 앱에서 같은지 함께 확인할 것」 항목} \\
\midrule
반사 손실 & 0.5 dB/회 & \hl{0.5 dB/회} & 0.5 dB/회 &
  RLS-19 D\textsubscript{RV} 0 dB 이상 임의 값(표 8). ISO 전파는 반사계수 ρ(0.2 초과 1 이하), RLS-90 전파는 D\textsubscript{E} & \ok \\
반사 거리 제한 & 반사면: 수음점 200 m, 음원 50 m 안 & 반사면–수음점 \hl{200 m}, 반사면–음원 \hl{50 m}(경로의 모든 반사면) &
  반사면–수음점 100 m, 음원–수음점 간선 200 m·골목 50 m, 2차 반사면 간격 50 m, 반사면–음원 제한 없음 &
  네 가지 거리를 각각 임의 값(m)으로 지정. 반사면–음원 제한은 끌 수 있음 & \ok \\
차폐량 상한 & 1000 dB (상한 없음) & \hl{상한 없음} & RLS-19 전파: 상한 없음. ISO 전파: 단일 20 dB·다중 25 dB &
  RLS-19 전파는 식 (15) 원문대로 상한 없음(고정). ISO 전파는 상한 값 지정 가능 & \ok \\
평가값 & 주간 LrD (06–22시) & \hl{주간 L\textsubscript{r,T}}(= LrD) & 주간 L\textsubscript{r,T}·야간 L\textsubscript{r,N} 함께 계산 &
  시간대 고정(주간 06–22시, 야간 22–06시) & \ok \\
\bottomrule
\end{tabularx}}

\vspace{1mm}
{\footnotesize 4절 표의 SORI 열을 위 결과로 바꾸면: 반사 차수 \eq{2차} · 노면 보정 \eq{0 dB} · 신호교차로 보정 \eq{없음} · 다중반사 \eq{없음} ·
건물 높이 \eq{층수 × 3 m} · 측면 회절 \eq{없음} · 기준 표기 \eq{RLS-19}.}

%% ------------------------------------------------------------------
\sect{2. 나머지 계산 조건(「SoundPLAN 계산 조건 정리」 1절 설정표)}

{\footnotesize
\renewcommand{\arraystretch}{1.25}
\begin{tabularx}{\linewidth}{@{}>{\bfseries}P{19mm}P{24mm}P{30mm}P{28mm}Lc@{}}
\toprule
\rowcolor{hdr} 항목 & SoundPLAN & SORI 계산 조건 & SORI 기본값 & SORI 변경 가능 범위 & 판정 \\
\midrule
도로 구간 분할 & 거리/직경 계수 8, 최소 1 m & \hl{1 m}(수음점 거리 s ≥ 8 m에서 l ≤ s/8) & 간선 3 m, 골목 5 m &
  간선·골목 따로 임의 길이(m). RLS-19 3.2 주석 l ≤ s/2 조건 진단 기록 & \ok \\
탐색 반경 & 5,000 m, 허용 오차 0.1 dB & 1,000 m(음원이 모두 600 m 안 → 5 km와 같은 결과) & 간선 300 m, 골목 150 m &
  간선·골목 따로 임의 거리(m) & \ok \\
교통량 & 현장 조사값, 화물차 모두 Trucks1(p1) & SoundPLAN 입력값: 동행 1,046 승용 + 388 Trucks1, 서행 838.2 승용 + 267.9 Trucks1(대/h), p2 = 0 &
  교통량 CSV. 자료 없는 도로는 차로 수별 가정 교통량 &
  도로(UFID·도로명)마다 주간·야간 시간당 교통량, p1·p2, 방향별 교통량·화물차 비율 & \ok \\
속도 & 남측 50, 북측 30 km/h (차종 공통) & 남측(동행) 50, 북측(서행) 30 km/h & 자료에 없으면 간선 50, 골목 30 km/h &
  방향별·차종별(승용·화물). 30 km/h 미만은 30 km/h(RLS-19 2장) & \ok \\
음원선 & 차로마다 & \hl{차로마다}, 사가정로 3 + 2 차로 & 주행 방향마다 1개선(RLS-19 Bild 6) &
  방향마다 / 차로마다. 방향별 차로 수 도로별 지정 & \ok \\
도로 높이·경사 & 지형 높이, 경사 한계 100 \%, 평활 0 m & \hl{지형 그대로}, 평활 0 m & 차도 면 평탄화(종단 20 m), 지형 평활 5 m &
  지형 그대로 / 평탄화, 평활 길이 임의. 경사 보정은 RLS-19 식 (7) & \ok \\
지형 & 등고선·표고점 & \hl{등고선·표고점만} & 등고선·표고점(DEM 없음) &
  DEM 래스터를 보조 또는 주 자료로 추가 & \ok \\
기타 객체 & 방음벽·감쇠 영역·지면 영역 없음 & \hl{사용 안 함} & 사용 안 함 &
  방음벽 CSV(반사·흡음), 녹지 지면 폴리곤 & \ok \\
수음점 & M1·M2, 높이 1.31 m & SoundPLAN과 \hl{같은 좌표·높이} & 격자·건물 외벽 켬, 점 수음점 없음 &
  점 수음점 CSV(좌표·높이), 격자, 건물 외벽(층별) & \ok \\
격자 & 5 m 간격, 지면 위 2 m & SoundPLAN 격자점과 \hl{같은 좌표}, \hl{2 m} & 5 m 간격, 지면 위 4 m &
  간격·높이(여러 높이)·범위 임의 & \ok \\
격자 보간 & 9×9, 차이 0.15 dB, 40 dB 미만 항상 보간 & 모든 점 직접 계산 & 모든 점 직접 계산 &
  보간 기능 없음(그림용 세분 격자는 보간이 아니라 더 촘촘히 계산) & \nb{맞출 수 없음} \\
\bottomrule
\end{tabularx}}

%% ------------------------------------------------------------------
\sect{3. SORI 계산 기능 한눈에}
SORI는 국토지리정보원 수치지형도(건물·도로·등고선·표고점)와 교통량 CSV를 입력으로 도로교통소음 지도를 계산한다. 설정은 YAML 파일 하나에 모이며,
지정하지 않은 항목은 기본값을 쓴다. 아래는 SoundPLAN 비교와 관련된 기능만 추렸다.

{\small
\renewcommand{\arraystretch}{1.2}
\begin{tabularx}{\linewidth}{@{}>{\bfseries}p{26mm}Lp{42mm}@{}}
\toprule
\rowcolor{hdr} 구분 & 기능 · 설정 범위 & 설정 키 \\
\midrule
계산 기준 & 전파 방법 3가지: RLS-19(배출·전파) / RLS-19 배출 + ISO 9613-2 전파 / RLS-90(배출·전파). RLS-19·RLS-90은 A 가중 단일값, ISO 9613-2는 옥타브 대역 계산 가능 & \f{propagation\_method} \\
반사 & 차수 \hl{0 \textasciitilde{} 6차}(격자·외벽·점 따로). 반사 손실 D\textsubscript{RV}(RLS-19)·반사계수 ρ(ISO, 0.2 초과 1 이하)·D\textsubscript{E}(RLS-90). 반사면–수음점·반사면–음원·음원–수음점 거리 제한과 2차 반사면 간격. 수음점이 놓인 외벽 반사 제외 규칙(마지막 반사만 / 모든 반사) & \f{propagation.reflection.*} \\
회절 & 건물·지형 위 회절(단일·다중). 측면 회절은 ISO 9613-2 전파에서만(켬/끔, 감싸 돌기 0 \textasciitilde{} 20회). 차폐 상한: RLS-19 없음, ISO 20 / 25 dB(설정 가능) & \f{propagation.lateral}, \f{propagation.barrier} \\
배출 보정 & 노면 13종(RLS-19 표 4a·4b), 종단경사(식 7), 신호·회전 교차로 D\textsubscript{K}(표 5, 자동 판정 또는 CSV), 다중반사 D\textsubscript{refl}(식 8, 켬/끔) & \f{input.surface\_map}, \f{input.junctions}, \f{model.canyon} \\
교통량·음원 & 도로별 주간·야간 교통량, p1·p2, 방향별 비율·속도. 교통량 자료 없는 도로의 차로수별 가정값. 음원선 방향별/차로별, 차로폭 자동/고정, 분할 길이 & \f{input.traffic}, \f{model.source\_lines}, \f{model.segment\_*} \\
지형·건물 & 등고선·표고점 지형(+ DEM 선택), 지형 평활, 차도 면 평탄화/그대로, 건물 층수 × 층고 또는 높이 필드, 방음벽, 녹지 지면 & \f{model.*}, \f{input.dem}, \f{input.barriers} \\
수음점 & 격자(간격·높이·범위, 경계 세분), 건물 외벽(층별·조각 길이), 점 수음점(좌표·높이) & \f{receivers.*} \\
결과 & 주간 L\textsubscript{r,T}·야간 L\textsubscript{r,N}, 격자·외벽·점 CSV와 지도 그림, 점별 도로 기여도, 감쇠항별 중간값(D\textsubscript{div}, D\textsubscript{atm}, D\textsubscript{gr}, D\textsubscript{z}) & \f{receivers.points.}\newline\f{road\_contrib},\newline\f{propagation.terms\_csv} \\
\bottomrule
\end{tabularx}}

\vspace{1mm}
{\footnotesize SORI 기본값은 이번 비교 설정과 다르다(전파 ISO 9613-2, 반사 3차, 측면 회절 켬, D\textsubscript{K}·D\textsubscript{refl} 켬, 차도 면 평탄화).
SoundPLAN 비교에는 1·2절의 값으로 바꾼 설정 파일을 따로 썼다.}

%% ------------------------------------------------------------------
\sect{4. 같게 맞출 수 없거나 해석이 들어간 항목}
\begin{enumerate}[leftmargin=1.4em]
\item \hl{격자 보간(맞출 수 없음).} SoundPLAN은 9×9 칸 모서리를 먼저 계산하고 조건이 맞으면 나머지 점을 보간한다(차이 0.15 dB 이내, 40 dB 미만 구역은 항상 보간).
  SORI는 모든 점을 직접 계산한다. 따라서 SoundPLAN 격자값 일부는 직접 계산값이 아니며, 점별 비교에서 0.15 dB 수준 차이는 의미가 없다.
  수음점 M1·M2는 보간과 관계없다.
\item \hl{반사면–음원 50 m 판정(해석).} SoundPLAN 설정 「max.\ reflection distance to source」 50 m를 SORI에서는
  「경로의 \hl{모든} 반사면이 음원(분할 구간)에서 50 m 안, 수음점에서 200 m 안」으로 구현했다(\f{r\_refl\_source: 50}, \f{r\_facade: 200}).
  SoundPLAN이 이 거리를 분할 구간마다 재는지, 도로 전체 기준으로 재는지, 첫 반사면에만 적용하는지는 설정 화면만으로 확인할 수 없다.
\item \hl{탐색 반경 1 km.} SoundPLAN 5 km와 값은 다르지만, 교통량이 있는 음원(사가정로)이 모두 수음점에서 600 m 안에 있어 계산 결과는 같다.
\item \hl{건물 안 격자점.} SoundPLAN 격자 중 수치지도 건물 안에 들어가는 약 150점은 원인을 확인하기 전까지 비교에서 뺀다(「SoundPLAN 계산 조건 정리」 3절 권고와 같음).
\end{enumerate}

%% ------------------------------------------------------------------
\clearpage
\sect{부록. SORI 설정 파일(이번 계산)}
\f{validation/soundplan/sp\_compare\_1109/sp\_settings\_1117/sp\_refl\_1119/config.yaml}에서 계산 조건 부분만 옮겼다(주석은 줄임).
{\footnotesize
\begin{verbatim}
propagation_method: rls19              # 배출·전파 모두 RLS-19 (측면 회절 적용 안 함)
input:
  map: {folders: [전농사거리_수치지도]}  # 등고선·표고점·도로·건물, DEM 없음
  traffic: {path: traffic_soundplan_1119.csv}   # SoundPLAN 입력 교통량, 화물차 전부 Lkw1
  surface_map: {"아스팔트": "ref", "아스팔트콘크리트": "ref"}   # 기준 노면, D_SD = 0
  junctions: {auto_signal_min_lanes: 99}       # 신호교차로 보정 없음, D_K = 0
model:
  floor_height: 3.0                    # 층수 × 3 m
  source_lines: lane                   # 차로마다 음원선
  canyon: {enabled: false}             # D_refl = 0
  segment_main: 1.0                    # 분할 1 m
  road_surface: dtm                    # 도로 높이 = 지형
  terrain_smooth_m: 0.0                # 평활 0 m
receivers:
  points: {path: points_sp.csv}        # SoundPLAN 격자점(2 m)·수음점 M1·M2(1.31 m)
propagation:
  max_dist_main: 1000.0
  reflection: {D_RV: 0.5,
               max_order_grid: 2, max_order_facade: 2, max_order_point: 2,
               r_facade: 200.0, r_refl_source: 50.0,
               r_source_main: 400.0, r_source_minor: 400.0, r_second: 520.0}
\end{verbatim}}

\vfill
{\footnotesize\color{gray} 근거: 「SoundPLAN 계산 조건 정리」(2026-11-17, SoundPLAN 9.1 설정 화면) · SORI v0.9 \f{sori/config.py}, \f{sori/rls19.py}.
본 문서는 Claude(Anthropic의 AI)를 활용하여 작성하였음.}

\end{document}
